\documentclass[10pt]{article}
\usepackage[a4paper,margin=27mm]{geometry}
\usepackage{amsmath,amssymb,amsthm}
\usepackage[T1]{fontenc}
\usepackage{lmodern,microtype}
\usepackage[colorlinks=true,allcolors=blue]{hyperref}
\hypersetup{pdftitle={Ambient Dependence of Extremal-Exchangeable Independence},
 pdfauthor={Alec Kriebel},
 pdfsubject={Unrefereed AI-assisted mathematical research note},
 pdfkeywords={exchangeability, independence, free products, extreme states}}
\setlength{\emergencystretch}{2em}
\newtheorem{theorem}{Theorem}
\newtheorem{proposition}[theorem]{Proposition}
\newtheorem{lemma}[theorem]{Lemma}
\newcommand{\C}{\mathbb C}
\newcommand{\F}{\mathcal F}
\newcommand{\E}{\mathcal E}
\newcommand{\id}{\mathrm{id}}
\title{Ambient Dependence of Extremal-Exchangeable Independence}
\author{Alec Kriebel\\
  \small ORCID: \href{https://orcid.org/0009-0001-9320-500X}{0009-0001-9320-500X}}
\date{12 October 2026}
\begin{document}
\maketitle
\begin{abstract}
Liebscher asked whether independence defined by an extreme exchangeable
free-product state needs to retain its ambient algebra. We give a
state-preserving inclusion of a four-point algebra into \(M_4(\C)\) for
which two distinct proper coordinate subalgebras, intersecting only in
the scalars, are dependent downstairs and independent upstairs.
Positivity forces every downstairs witness to be the same nontrivial
character mixture; upstairs a pure vector state supplies a pure folding
witness. Pure folding is an established construction, credited below.
We also exhibit an extreme exchangeable algebraic state with no positive
extension to a larger algebra, addressing the source's separate
non-lifting concern. In contrast, extreme invariant states do lift
along equivariant unital \(C^*\)-inclusions. The definition is unrestricted:
extremality is among all exchangeable states, without a fixed-marginal,
faithfulness or trace condition.
\end{abstract}

\section{The original question and convention}
Let \(A\) be a unital complex *-algebra and let \(\varphi\) be a positive
unital linear functional. Write
\(\F(A)=*_{i\in\mathbb N}A\) for the unital algebraic free product,
with injections \(j_i\), and let \(\E(A)\) be its convex set of states
invariant under finite permutations of the copies. In the \(C^*\)-category
we use the universal unital \(C^*\)-free product and continuous states.

For unital *-subalgebras \(B,C\subseteq A\), Liebscher's definition
on printed page 540 of \cite{Liebscher} declares them independent
when some \(\omega\in\operatorname{ext}\E(A)\) satisfies
\begin{equation}\label{law}
 \omega\circ\bigl((j_1|_B)*(j_2|_C)\bigr)
 =\varphi\circ(\iota_B*\iota_C)
 \quad\hbox{as states on }B*C.
\end{equation}
Here \(\iota_B,\iota_C\) are the original inclusions, and the maps
from \(B*C\) are supplied by the universal property. Equality means
the entire joint law. Extremality is in \(\E(A)\), not in a slice
with prescribed marginals. Item (1) asks whether retaining the ambient
algebra is necessary, and motivates this by a possible failure of
extreme-state lifting. We distinguish those two issues.

\section{A finite-dimensional counterexample}
\begin{theorem}\label{main}
Put \(A=\C^{\{00,01,10,11\}}\), let \(p(s,t)=s,\ q(s,t)=t\), and set
\[
 B=\operatorname{span}\{1,p\},\qquad
 C=\operatorname{span}\{1,q\},\qquad
 \varphi=\tfrac12\delta_{00}+\tfrac12\delta_{11}.
\]
Embed \(A\) diagonally in \(M_4(\C)\), in the displayed point order, and let
\[
 u=(e_{00}+e_{11})/\sqrt2,\qquad \varphi'(T)=u^*Tu.
\]
This inclusion preserves the state. The subalgebras \(B,C\) are distinct,
proper and have scalar intersection. They are not independent in
\((A,\varphi)\), but are independent in \((M_4(\C),\varphi')\)
under \eqref{law}, in both free-product categories.
\end{theorem}

\begin{proof}
The functions \(1,p,q,pq\) form a basis of \(A\).
Evaluating \(a+bp=c+dq\) at \(00,10,01\) gives \(a=c,\ b=d=0\).
The stated subalgebra properties follow. Restriction of the vector
state to the diagonal gives exactly \(\varphi\).

Suppose that an exchangeable state \(\omega\) satisfies \eqref{law}
downstairs. In its GNS representation put \(P_i=\pi(j_i(p))\),
\(Q_i=\pi(j_i(q))\), with cyclic vector \(\Omega\).
The joint law and exchangeability imply, for every \(i\ne j\),
\[
 \langle\Omega,P_i\Omega\rangle
 =\langle\Omega,Q_j\Omega\rangle
 =\langle\Omega,P_iQ_j\Omega\rangle
 =\langle\Omega,Q_jP_i\Omega\rangle=\tfrac12 .
\]
Consequently
\[
 \|(P_i-Q_j)\Omega\|^2=0\qquad(i\ne j).
\]
Choosing a third index shows that all \(P_i\Omega\) and \(Q_i\Omega\)
are the same vector \(v\). Each projection \(R\) in this family satisfies
\(R\Omega=v\) and \(Rv=R^2\Omega=v\). Thus every nonempty finite word
in the generators sends \(\Omega\) to \(v\) and has moment \(1/2\).
These words and the unit span the algebraic free product.

Let \(\chi_0,\chi_1\) evaluate every copy at \(00,11\), respectively.
The preceding calculation proves, for every possible witness,
\[
                  \omega=\tfrac12\chi_0+\tfrac12\chi_1.
\]
The two characters are distinct exchangeable states. Therefore
\(\omega\) is not extreme. Equality on dense algebraic words gives
the same conclusion for the universal \(C^*\)-free product.
This excludes every downstairs witness, rather than just a chosen one.

Upstairs, let \(f:\F(M_4)\to M_4\) fold all copies identically.
The vector state \(\varphi'\) is pure, and its pullback
\(\omega'=\varphi'\circ f\) is pure by Lemma~\ref{pure}.
The folding map is unchanged by permutations, so \(\omega'\) is
exchangeable and hence extreme in the exchangeable state space.
For every word from \(B*C\), folding multiplies the original
diagonal matrices in their original order. Restriction of
\(\varphi'\) to \(A\) proves \eqref{law} on the entire algebra.
Moreover \(\omega'\circ j_i=\varphi'\) on every full copy.
\end{proof}

The algebraic GNS calculation above needs only Cauchy--Schwarz.
For clarity, its null space is a left ideal: if \(\eta(x^*x)=0\),
Cauchy--Schwarz applied to \(x\) and \(a^*ax\) gives
\(\eta(x^*a^*ax)=0\), so \(ax\) is null. Projection multiplication
therefore acts on the quotient domain. No null-vector identity was
silently promoted to an operator identity.

\begin{lemma}\label{pure}
A pure state pulls back to a pure state along a surjective unital
*-homomorphism. A unit vector state on \(M_d(\C)\) is pure.
Both statements hold for positive algebraic states.
\end{lemma}
\begin{proof}
In a convex decomposition of \(\eta\circ f\), positivity forces each
constituent to vanish on \(x^*x\) for \(x\in\ker f\).
Cauchy--Schwarz then makes it vanish on \(x\). Each constituent
descends to a unital positive functional on the quotient: lift
any \(a\) to \(z\), and evaluate \(a^*a\) as \(z^*z\).
Purity of \(\eta\) makes the decomposition trivial.
For a vector state, let \(e=uu^*\). The zero expectation of \(1-e\)
in every constituent kills all blocks outside \(e\), by
Cauchy--Schwarz. Since \(eTe=(u^*Tu)e\) and each constituent
has value one on \(e\), each is the original vector state.
\end{proof}

Pure folding is the known maximal-amalgamation construction:
Dykema, K\"ostler and Williams \cite[Proposition 8.7]{DKW} characterize
its quantum symmetric extremality, and their Theorem 8.1 relates
this to ordinary symmetric extremality. Lemma~\ref{pure} gives the
needed special case directly, with the stronger conclusion of purity
among all states.

The half weight is inessential. Replacing it by
\((1-t)\delta_{00}+t\delta_{11}\), \(0<t<1\), and using
\(\sqrt{1-t}\,e_{00}+\sqrt t\,e_{11}\), gives the same proof with
moment \(t\). At \(t=0,1\) the downstairs mixture becomes a character.
If coincident subalgebras are allowed, the smaller example
\(A=B=C=\C^2\subset M_2\) works with the faithful downstairs fair
state and its coherent pure extension.

\section{The literal lifting concern}
\begin{theorem}\label{nonlift}
There is a unital inclusion of algebraic *-algebras for which an
extreme exchangeable state on the smaller free product has no
positive extension to the larger one.
\end{theorem}
\begin{proof}
Take \(R=\C[x]\subset S=\C[x,x^{-1}]\), with both \(x\) and \(x^{-1}\)
self-adjoint. Evaluation at zero in every copy gives a character
\(\chi\) on \(\F(R)\), hence a pure exchangeable state.
A positive unital extension \(\eta\) to \(\F(S)\) would satisfy,
in the first copy, \(X^*=X,\ Y^*=Y,\ XY=1,\ \eta(X^2)=0\).
Cauchy--Schwarz gives the contradiction
\[
          1=|\eta(XY)|^2\le \eta(X^2)\eta(Y^2)=0.
\]
\end{proof}
Both \(R,S\) admit the common state given by evaluation at one.
The forbidden \(\chi\) is a different state; the assertion concerns
the unrestricted extreme-state sets in the source, not failure
to extend that specified common marginal.

\begin{proposition}\label{Clift}
For an equivariant unital \(C^*\)-inclusion \(D\subset E\) under the group
of finitely supported permutations of \(\mathbb N\), every extreme invariant
state of \(D\) has an extreme invariant extension to \(E\).
\end{proposition}
\begin{proof}
Extend the state to \(E\) by the \(C^*\)-state extension theorem.
Average the extension over \(S_n\) and take a weak-* cluster point,
using a subnet when necessary. Every fixed finite permutation
fixes all sufficiently late averages; the limit is an invariant
extension. Its compact convex extension fiber is a face of the
invariant state space of \(E\), since restriction of a decomposition
would decompose the extreme state of \(D\). An extreme point of
this nonempty compact face is extreme in the whole space.
\end{proof}
The canonical map between universal unital free products of a
\(C^*\)-inclusion is injective. One may check this by extending each
copy representation using state extension and cyclic GNS
dilations, padding them to a common Hilbert space, and applying
the universal property. The verification package gives this
argument in detail. Thus Proposition~\ref{Clift} applies to those
products. Extreme states can lift while their restrictions fail
to remain extreme. Theorem~\ref{main} proves ambient dependence
in that opposite direction; Theorem~\ref{nonlift} concerns the
algebraic category, where \(C^*\)-extension and compactness are unavailable.

\section{Scope, attribution and verification}
The printed free-product definition is the target. We impose no
global commutativity on its state algebra merely because \(A\)
is commutative. A classical commutative/tensor interpretation
requires a different coproduct or states factoring through a
commutative quotient; it is not assumed here.

The counterexample uses nonfaithful states, as permitted by the
definition, and no tracial witness. It does not address faithful-only,
tracial-only, trivial-tail or separately specified normal W*-definitions.
Adding a prescribed full marginal while retaining extremality in
\(\E(A)\) leaves the upstairs witness valid. Changing the convex
set to the fixed-marginal slice changes the result: the forced
downstairs mixture is extreme in that slice, as independently proved
in the supporting audit.

The known pure-folding mechanism is credited above. Our source audit
does not certify historical-first priority, and we make no such claim.
The note supplies explicit counterexamples and full deductions for
the stated original question, rather than resolving all the other
questions in the workshop contribution.

The source archive contains the manuscript, immutable authored input,
independent proofs and rational verification scripts, source hashes
and the priority audit. Finite checks supplement the all-state,
all-word and infinite-copy proofs; they are not formal certification.
AI tools were used extensively in proof development, reproduction,
adversarial review and preparation. All listed reviews were AI-assisted.
This preprint is unrefereed and has not received human peer review.

\begin{thebibliography}{9}
\bibitem{Liebscher}
V.~Liebscher, On Quantum Independence, pp.~540--541 in
B.~V.~R.~Bhat, U.~Franz and M.~Skeide (organizers),
\emph{Mini-Workshop: Product Systems and Independence in Quantum Dynamics},
Oberwolfach Reports \textbf{6} (2009), 493--548.
\href{https://doi.org/10.4171/OWR/2009/09}{doi:10.4171/OWR/2009/09}.
\bibitem{DKW}
K.~J.~Dykema, C.~K\"ostler and J.~D.~Williams,
Quantum symmetric states on free product \(C^*\)-algebras,
\emph{Trans. Amer. Math. Soc.} \textbf{369} (2017), 645--679.
\href{https://doi.org/10.1090/tran6661}{doi:10.1090/tran6661};
\href{https://arxiv.org/abs/1305.7293v3}{arXiv:1305.7293v3}.
\end{thebibliography}
\smallskip\noindent\small
Version 1.0. Author: Alec Kriebel. License: CC BY 4.0.
\end{document}
